\section{Prefix Controls}
There are 3 or 4 letter prefixs that allow special control of your geometry and simulation. These help with assigning porous media regions, Multiple Reference Frame (MRF) regions, or even more control over meshing for certain geometries.

\subsection{Rotating Geometries}
The \codeword{ROTA} prefix allows the assignment of rotating wall to geometries such as wheels. A usage example is shown below.
\begin{lstlisting}
    [DEFAULTS]
    DEFAULT_MODULES = defaultGlobalControl defaultGlobalSimControl defaultBCSetup defaultGlobalDecomposition defaultGlobalMaterial defaultGlobalRefinement defaultPost
    ADDON_MODULES =

    [TITLES]
    CASENAME = casename
    CASEDESCRP = default description
    JOBCODE = TS

    [GEOMETRY]
    GEOM = body.obj.gz,1,8,5,1.1,high
            ROTA-fr-wh-lhs,1,8,5,1.1,high
            ROTA-fr-wh-rhs,1,8,5,1.1,high
            ROTA-rr-wh-lhs,1,8,5,1.1,high
            ROTA-rr-wh-rhs,1,8,5,1.1,high
\end{lstlisting}
Notice the prefix \codeword{ROTA}, it lets \codeword{caseSetup} know that this geometry will have a rotating wall boundary condition applied. Upon running \codeword{caseSetup}, if it's the first time with this set of geometry, \codeword{caseSetup} will stop and let you know that the following blocks have been added. The user should check that these are the appropriate settings for each wheel geometry.

\textbf{NOTE:} If using a plinth, ensure that the plinth is a separate geometry, else the automatic wheel center calculation will include the plinth in the radius calculation resulting in an incorrect radius and angular velocity!

\begin{lstlisting}
    [ROTA-fr-wh-lhs]
    WH_RAD = default
    WH_CENTER = default
    WH_AXIS = default
    ROT_WH = True

    [ROTA-fr-wh-rhs]
    WH_RAD = default
    WH_CENTER = default
    WH_AXIS = default
    ROT_WH = True

    [ROTA-rr-wh-lhs]
    WH_RAD = default
    WH_CENTER = default
    WH_AXIS = default
    ROT_WH = True

    [ROTA-rr-wh-rhs]
    WH_RAD = default
    WH_CENTER = default
    WH_AXIS = default
    ROT_WH = True

\end{lstlisting}
\begin{itemize}
    \item \codeword{WH_RAD} [[\codeword{float,m}]] - wheel radius, distance from wheel center to ground (\codeword{r}, \codeword{default})
    \item \codeword{WH_CENTER} [[\codeword{float,m}]] [[\codeword{float,m}]] [[\codeword{float,m}]] - wheel center coordinates in global coordinates system (\codeword{x} \codeword{y} \codeword{z}, \codeword{default})
    \item \codeword{WH_AXIS} [[\codeword{float,m}]] [[\codeword{float,m}]] [[\codeword{float,m}]] - wheel axis of rotation from, using right hand rule with respect to the global coordinate system (\codeword{x} \codeword{y} \codeword{z}, \codeword{default})
    \item \codeword{WH_AXIS} [\codeword{bool}] - activate rotating wall or use no-slip condition (\codeword{true}, \codeword{false})
\end{itemize}

\subsection{Fan Pressure-Jump Condition}
The prefix \codeword{FAN} defines a geometry-driven circular fan disk in a snappyHexMesh case. The disk is meshed as a baffle/face zone and is assigned the ESI OpenFOAM \codeword{fan} pressure-jump condition. This is intended for applications such as an electric fan mounted upstream of a radiator. It is independent of the \codeword{POR} radiator model: a radiator may be represented by a porous volume, while the fan is represented by a separate planar disk.

The FAN geometry must be a single, approximately planar circular surface in an ASCII OBJ or STL file. It should not contain the fan frame, motor, blades, radiator, or other solids. The geometry file is supplied in the normal \codeword{[GEOMETRY]} section, using the \codeword{FAN-} prefix. FAN geometry is currently handled by the snappyHexMesh path; ANSA fan meshing is not enabled.

For example, the following setup places a fan disk upstream of a porous radiator. With \codeword{default} values, the center and equivalent circular diameter are calculated from the disk, and the disk normal is oriented toward the target radiator geometry.

\begin{lstlisting}
    [GEOMETRY]
    GEOM = body.obj.gz,1,8,5,1.1,high
            POR-radiator.obj.gz,1,10,6,1.4,high
            FAN-radiator-fan.obj.gz,1,10,0,1.0,high

    [FAN-radiator-fan]
    FAN_MODEL = fan
    FAN_CENTER = default
    FAN_DIRECTION = default
    FAN_TARGET_GEOMETRY = POR-radiator
    FAN_DIAMETER = default
    FAN_PRESSURE_RISE = 250
    SAMPLE_FAN = True
\end{lstlisting}

The FAN configuration block is exposed automatically the first time the geometry is encountered. The available settings are:
\begin{itemize}
    \item \codeword{FAN_MODEL} [\codeword{str}] - ESI fan model selector. Use \codeword{fan} (\codeword{pressureJump} is accepted as an alias).
    \item \codeword{FAN_CENTER} [\codeword{float,m}] - disk center in global coordinates, or \codeword{default} to use the area-weighted surface centroid.
    \item \codeword{FAN_DIRECTION} [\codeword{float}] - three-component direction vector, or \codeword{default} to orient the disk normal toward \codeword{FAN_TARGET_GEOMETRY}. Explicit direction takes precedence over the target geometry.
    \item \codeword{FAN_TARGET_GEOMETRY} [\codeword{str}] - geometry name without its extension. Required when \codeword{FAN_DIRECTION = default}; for example, \codeword{POR-radiator}.
    \item \codeword{FAN_DIAMETER} [\codeword{float,m}] - disk diameter, or \codeword{default} to calculate the equivalent circular diameter from surface area.
    \item \codeword{FAN_PRESSURE_RISE} [\codeword{float,Pa}] - constant pressure jump applied by the ESI fan condition. The example applies a 250 Pa pressure rise. Pressure curves are not currently accepted by the \codeword{caseSetup} configuration; only one scalar value is supported.
    \item \codeword{SAMPLE_FAN} [\codeword{bool}] - fan sampling switch retained in the setup metadata; FAN reporting is not yet implemented.
\end{itemize}

The generated pressure boundary uses the ESI syntax \codeword{type fan}, \codeword{patchType cyclic}, and a constant \codeword{jumpTable}. The pressure-rise sign follows the resolved FAN direction. If the fan accelerates flow in the opposite direction, provide an explicit \codeword{FAN_DIRECTION} or reverse the target geometry arrangement before running the case. Verify the generated boundary condition with the installed ESI OpenFOAM release (\codeword{of2206} or \codeword{of2606}) before production use.

Using \codeword{default} option will trigger the automatic calculation but manual input of values is still possible. To speed up \codeword{caseSetup} it is recommended to copy the values output by the calculation in the termial and input them into the caseSetup if coordinates and axis values are not expected to change for subsequent runs.


\subsection{Porous Media}
When radiators are used in a simulation, it would be very computationally expensive to model each radiator fin. So instead we use the Darcy-Forcheimer porous media calculation to model the pressure drop across a radiator which allows us to estimate the mass flow across the radiator and resulting forces.

The prefix \codeword{POR} is used to define the porous media geometry. \textbf{CAUTION:} The porous media geometry must only contain the "box" or region that represents the radiator fin region, do not include the frame or any other part of the radiator. The box which represents the fin region should be a simple prism. The inlet and outlets should have their own PID with the keyword \codeword{-inlet} or \codeword{-outlet}. There should only be 1 of each PID!

The mass flow calculations are performed on each pid in the porous media, therefore including the walls or any other surfaces where flow isn't moving \textit{across} would make the calculations erronous.

Including the prefix \codeword{POR} on a geometry would expose the following block:
\begin{lstlisting}
    GEOM = trial005-rad.obj.gz,1,8,6,1.4,high
            fw-H05A00.obj.gz,1,8,6,1.4,high
            rw-R11I01.obj.gz,1,8,6,1.4,high
            POR-rad-001.obj.gz,1,10,6,1.4,high
            ROTA-fr-wh-lhs-v3-3.obj.gz,1,8,6,1.4,high
            ROTA-fr-wh-rhs-v3-3.obj.gz,1,8,6,1.4,high
            ROTA-rr-wh-lhs-v3-3.obj.gz,1,8,6,1.4,high
            ROTA-rr-wh-rhs-v3-3.obj.gz,1,8,6,1.4,high

    [POR-rad-001]
    DCOEFFS = 1000
    FCOEFFS = 1000
    VEC1 = 0.7071 0 -0.7071
    VEC2 = 0 1 0
    POINT = 1.7875 -0.000583 0.342739
    SAMPLE_POR = True
\end{lstlisting}
the name of the block will be the name of the geometry just without the extension to identify it.
\begin{itemize}
    \item \codeword{DCOEFFS} [\codeword{float}] - Darcy coefficient attained from experimental testing of pressure drop vs flow rate of the radiator (\codeword{float})
    
    \item \codeword{FCOEFFS} [\codeword{float}] - Forcheimer coefficient attained from experimental testing of pressure drop vs flow rate of the radiator (\codeword{float})
    
    \item \codeword{VEC1} [[\codeword{float,m}] [\codeword{float,m}] [\codeword{float,m}]] - vector defining the normal of the radiator exit (\codeword{float} \codeword{float} \codeword{float})
    
    \item \codeword{VEC2} [[\codeword{float,m}] [\codeword{float,m}] [\codeword{float,m}]] - vector in plane with the radiator exit  (\codeword{float} \codeword{float} \codeword{float})
    
    \item \codeword{POINT} [[\codeword{float,m}] [\codeword{float,m}] [\codeword{float,m}]] - point defining the "inside" of the radiator (\codeword{float} \codeword{float} \codeword{float})
    
    \item \codeword{SAMPLE_POR} [\codeword{str,m}] - to sample the mass flow through the radiator (\codeword{true}, \codeword{false})
    

\end{itemize} 
